\documentclass[11pt]{article}

% Apprentissage de bibliothèque défaisable — brouillon de soumission (FR). Compile : pdflatex (deux fois).
\usepackage[T1]{fontenc}
\usepackage[utf8]{inputenc}
\usepackage[margin=1in]{geometry}
\usepackage{booktabs}
\usepackage{amsmath,amssymb}
\usepackage{listings}
\usepackage{xcolor}
\usepackage[hidelinks]{hyperref}

\lstset{
  basicstyle=\ttfamily\footnotesize,
  breaklines=true,
  columns=fullflexible,
  frame=single,
  framesep=4pt,
  xleftmargin=2pt,
}

\newcommand{\K}{\textsf{K1}}
\renewcommand{\abstractname}{Résumé}
\renewcommand{\refname}{Références}
\title{\textbf{Apprentissage de bibliothèque défaisable :\\ des méthodes typées à contrats d'exécution qui se désapprennent à la dérive}}
\author{Nathanael Braun}
\date{2026 \\[2pt] {\normalsize\itshape Préprint v1 --- rapport de système / d'expérience}}

\begin{document}
\maketitle

\begin{abstract}
Les agents fondés sur de grands modèles de langage réutilisent leur travail passé via une mémoire floue ---
recherche (RAG), raisonnement à partir de cas (CBR), bibliothèques de compétences en prose. Toutes rappellent par
similarité de surface, et aucune ne sait représenter une \emph{prémisse devenue fausse} : quand le monde dérive sans
changer la requête, la réponse en cache reste le plus proche voisin, et elle est servie quand même. Les
bibliothèques apprises statiques ont le défaut inverse --- saines une fois apprises, mais incapables de
\emph{désapprendre}.

Nous présentons l'\textbf{apprentissage de bibliothèque défaisable} : une bibliothèque apprise de méthodes typées
et composables, chacune dotée d'un \textbf{contrat d'exécution défaisable}. La garantie d'une méthode est
\emph{supposée} à sa composition, \emph{vérifiée (assertée)} à son exécution, et \textbf{rétractée avec blâme}
quand elle échoue --- une étape de maintenance de la vérité --- après quoi la bibliothèque est \emph{révisée}
chirurgicalement, non jetée. La même structure typée qui rend une méthode canonicalisable rend aussi sa
réutilisation amortissable et sa composition vérifiable sur les seuls contrats, à contexte par appel borné. Aucun
mécanisme n'est nouveau en soi (JTMS, contrats à blâme, apprentissage de bibliothèque, logique de séparation) ; la
contribution est leur \textbf{composition} en une représentation unique où amortissement, vérification de
composition et \emph{désapprentissage principiel et sélectif à la dérive} coïncident --- ce qu'aucune des briques ne
fournit seule.

Nous évaluons chaque mécanisme isolément sur un moteur réel à base de règles (un simulateur déterministe plus un
modèle local en réel). Sous une dérive externe en cours de flux, les mémoires par rappel seul servent du périmé
tandis que le contrat déclaratif récupère l'exactitude \emph{sélectivement} et \emph{en sûreté de composition} ---
pour une méthode seule, à travers une chaîne de méthodes apprise, et en tête-à-tête face aux systèmes de mémoire
d'agents nommés (MemGPT, Reflexion, GraphRAG) --- borné par un plafond de canonicalisabilité mesuré que nous
appelons \K{}.
\end{abstract}

\section{Introduction}
Un agent qui résout de nombreux problèmes apparentés devrait devenir moins coûteux et plus fiable avec le temps.
La voie dominante consiste à mémoriser et rappeler : stocker des solutions ou compétences passées, retrouver la
plus proche, la réutiliser ou l'adapter. La génération augmentée par recherche~\cite{lewis2020}, le raisonnement à
partir de cas et les bibliothèques de compétences comme Voyager~\cite{wang2023} partagent cette forme et le même
angle mort : ils rappellent par similarité de \emph{surface} et ne représentent pas une \emph{prémisse devenue
fausse}. Quand le monde change d'une manière qui ne change \textbf{pas} la requête --- une réglementation se
durcit, un fait est audité et trouvé faux, une politique est révoquée --- la réponse en cache reste le plus proche
voisin, et elle est toujours servie. Les bibliothèques statiques (programmes appris, compétences
distillées~\cite{ellis2021,bowers2023}) ont le problème inverse : saines à l'apprentissage, elles ne savent pas
\emph{désapprendre}.

L'ingrédient manquant est un \textbf{contrat typé défaisable} sur chaque unité réutilisable. Empruntant aux
contrats logiciels avec blâme~\cite{findler2002} et à la vérification graduelle~\cite{bader2018}, une méthode
déclare ce qu'elle lit, écrit, exige et garantit, sur un alphabet de faits \emph{typé}. Pour une méthode
\emph{apprise} la garantie est une hypothèse induite : \textbf{supposée à la composition, vérifiée à l'exécution,
rétractée avec blâme} en cas de violation --- une opération de maintien de la vérité~\cite{doyle1979,dekleer1986}
--- après quoi la bibliothèque est révisée en \emph{spécialisant} la précondition fautive, non en la supprimant
(révision de théorie~\cite{ourston1994,richards1995} ; révision de croyances~\cite{agm1985}). La même structure
typée rend la réutilisation \textbf{amortissable} (une clé canonique stable) et la composition \textbf{vérifiable
sans ouvrir la boîte}~\cite{ohearn2001,reynolds2002}. Le tout est borné par une unique contrainte honnête, \K{} :
seule la structure typée, canonicalisable, amortit ; la prose véritable reste dans le modèle.

Nous employons \emph{apprentissage de bibliothèque} au sens de DreamCoder/Stitch~\cite{ellis2021,bowers2023}
(induire des méthodes typées réutilisables par abstraction~\cite{plotkin1970}), non un ajustement statistique ; on
pourrait parler d'\emph{induction de méthodes}. La nouveauté n'est \textbf{aucun} mécanisme isolé --- ni l'induction
ni la maintenance de la vérité, toutes deux d'un état de l'art vieux de plusieurs décennies --- mais leur
\emph{composition} : attacher un contrat défaisable à une méthode \emph{typée et composable}, c'est ce qui fait
tomber l'amortissement, la vérification de composition et le désapprentissage d'une seule représentation, là où
chaque brique seule n'en livre qu'un au plus. Le changement est de flot de contrôle : là où
une mémoire par similarité fait \texttt{requête $\to$ retrouver $\to$ réutiliser}, la nôtre fait
\texttt{requête $\to$ retrouver-le-contrat $\to$ vérifier $\to$ exécuter $\to$ contrôler $\to$ rétracter $\to$
spécialiser}. Le suffixe contrôler-et-rétracter est ce que les expériences isolent.

\paragraph{Contributions.} (1) Le cadrage des méthodes d'agent réutilisables comme \textbf{non-terminaux typés à
deux faces dotés d'un contrat d'exécution défaisable}, et la boucle supposer/vérifier/rétracter-blâmer/réviser. (2)
Une évaluation reproductible, isolant les mécanismes, sur un moteur réel --- la mémoire par rappel seul ne peut pas
désapprendre tandis qu'un contrat typé déclaratif récupère la dérive \emph{sélectivement, généralement,
sûrement-en-composition} (E2, simu + réel, avec une référence Invalidant équitable) ; transfert structurel sain que
le seul nombre d'appels ne certifie pas (E1) ; aucun faux-admis sur les compositions évaluées, chacune des trois
barrières ablée (E3) ; amortissement en gradient de la fraction canonicalisable (P4) --- en explicitant ce que
chaque expérience établit ou non.

\section{Approche}
\subsection{L'objet : une méthode à deux faces}
Une \textbf{méthode} est, pour son appelant, une boîte noire à contrat typé ; à l'intérieur, des
\emph{productions} la réalisent (séquence, branchement, map, fold). Formellement, un non-terminal de remplacement
d'hyperarêtes~\cite{habel1992,drewes1997} à sélection conditionnée par précondition~\cite{erol1994}. Nous tenons
deux régimes séparés par \emph{intention de conception} (nous ne prouvons pas la décidabilité ici). La
\textbf{grammaire des méthodes} est \emph{censée} rester décidable, via un rang de montage bien fondé et des
invariants de typage ; l'existence d'un plan HTN récursif étant indécidable en général~\cite{erol1996}, nous nous
restreignons au fragment bien fondé. L'\textbf{exécution} sur des données de taille d'exécution est l'inverse : une
couche explicitement bornée par budget, Turing-complète. Le lien à la logique monadique du second ordre~\cite{courcelle1990} est une motivation, non un
théorème établi ici.

\subsection{Le contrat : un triplet de séparation défaisable}
Une méthode déclare une empreinte de lecture, une empreinte d'écriture, une précondition, une postcondition et une
étiquette d'effet. La composition sous état partagé est le problème du cadre~\cite{mccarthy1969} ; nous le levons
par une discipline d'empreintes de logique de séparation~\cite{ohearn2001,reynolds2002} sur l'alphabet typé fini.
Pour une méthode \emph{apprise} la postcondition est une hypothèse induite : \textbf{supposée à la composition,
vérifiée à la stabilisation, rétractée avec blâme}~\cite{findler2002,bader2018}. Le moniteur d'exécution est le
JTMS~\cite{doyle1979}, donnant une sûreté \emph{éventuelle} (non statique) --- précisément le désapprentissage qui
manque aux références. Ce désapprentissage est symbolique (une rétractation de croyance par JTMS) et distinct du
\emph{machine unlearning} paramétrique~\cite{bourtoule2021}, qui efface l'influence de données d'entraînement des
poids d'un modèle.

\subsection{Pipeline et le plancher \K{}}
Une formulation est typée en un but ; une méthode est sélectionnée et composée sur contrats ; les cas la traversent ;
les traces distillent (anti-unification~\cite{plotkin1970} ; filtrées par MDL~\cite{ellis2021,bowers2023}) en de
nouvelles méthodes typées ; la dérive rétracte. Le repli universel est le \emph{plancher de micro-tâches} : tout ce
qui ne se réduit pas à une méthode typée en cache se réduit à une micro-tâche qu'un petit modèle traite aisément.
Un contrat \emph{manquant} coûte un appel modèle bon marché ; un contrat \emph{faux} est rattrapé par la
vérification d'exécution --- jamais une erreur silencieuse.

\section{Mise en œuvre}
Tous les mécanismes sont réalisés sur un moteur existant de graphe de faits typés à base de règles, cohérent par
JTMS, à concepts déclaratifs et stabilisation par chaînage avant, sans modification du cœur hormis une option de
requête additive. La clé de mémoïsation typée est le condensé de canonicalisation du moteur ; le vérificateur de
contrat défaisable (entailment à la composition sur domaines abstraits, assertion de postcondition à l'exécution,
trois barrières de sûreté) et la transformation de transfert structurel (relativiser-au-stockage /
lier-au-rejeu) sont des bibliothèques hôtes. Un exécuteur de cas durable à l'échelle est un artefact connu (un
réseau de workflow~\cite{vanderaalst1998} sur un magasin durable, dans la lignée d'AWS Step Functions Distributed
Map~\cite{aws2022}, du cache de contenu de Prefect, et de DBOS~\cite{dbos2022}) ; notre vue-croyance se situe
au-dessus. Le cycle de vie défaisable --- \textbf{supposer $\to$ vérifier $\to$ contrôler $\to$ rétracter $\to$
spécialiser} --- correspond un-à-un aux fonctions du moteur :

\begin{lstlisting}
select(but):                             # SUPPOSER (a la composition)
  M <- bibliotheque.match(but.faits_typ) #   cle typee ; un echec -> plancher de micro-taches
  supposer M.contrat                     #   checkCompose: post(prec) |= pre(M); escalade, jamais faux-admis

apply(M, cas):                           # VERIFIER + CONTROLER (a l'execution)
  cle <- digest(cas.premisse_typee)      #   cle canonique K1
  si memo.has(cle): retourner memo[cle]  #   amortir un cas type recurrent
  sortie <- run(M, cas)                  #   sinon deriver (appel modele / sous-graphe)
  si non assertPost(M.contrat, sortie):  #   la post tient ? + G1 cadre + G2 oracle-d'effet
    quarantaine(cas); blamer(M.contrat)  #   ne jamais valider une mauvaise sortie
    retourner
  memo[cle] <- sortie; retourner sortie

on ingest(fait):                         # RETRACTER + SPECIALISER (derive)
  pour e dans memo dependant de fait:    #   JTMS: re-verifier chaque post affectee
    si non satisfies(e.contrat.post, e.faits + {fait}):
      retracter(e); blamer(e.contrat)    #   desapprendre: evincer + attribuer le blame
  bibliotheque.reviser(blame): pre <- specialiser(pre)  # reviseOnBlame (CEGIS): restreindre, sans supprimer
\end{lstlisting}

\section{Expériences}
\subsection{Dispositif}
Les expériences utilisent les mécanismes réels du moteur à deux fidélités explicites. \textbf{E1 et E3 instancient
le moteur complet} (graphe + stabilisation + JTMS) ; \textbf{E2, P4 et E5 isolent les fonctions} réelles
(\texttt{digest}, \texttt{satisfies}, \texttt{canonValue}) depuis un banc, sans la boucle complète --- nous ne
prétendons pas qu'E2 exerce la rétractation JTMS native (il la ré-implémente sur le même prédicat réel). Le modèle
est un \textbf{simulateur déterministe} (oracle parfait de la règle \emph{courante} étant donné seulement ce que
révèle l'invite --- toute péremption/coût vient du mécanisme, non d'une erreur modèle) ou un \textbf{modèle local
réel} (Qwen3.6-27B (Q2\_K\_XL, MTP)). Chaque exécution est conditionnée par un \textbf{auto-test du banc} (sous le
simulateur le bras naïf doit être parfait, sinon avortement). Sept bras partagent une interface : Naïf,
Long-contexte, RAG, CBR (clé typée, sans re-vérification), Compétence (prose), \textbf{Invalidant} (cache à clé
typée + rappel grossier codé à la main : crochet d'invalidation, sans contrat typé), et \textbf{Struct} (la
bibliothèque typée au contrat défaisable). L'Invalidant sépare « possède une invalidation » de « possède un contrat
typé ».

\subsection{E2 --- défaisance à la dérive}
Un domaine d'approbation typé ($N{=}80$, deux classes auditées ; l'exécution réelle utilise $N{=}48$, une classe)
avec une invalidation de prémisse \emph{externe} en cours de flux : un audit de conformité marque une classe non
conforme, faisant basculer ses cas approuvés vers le refus. L'audit n'est \textbf{pas un champ d'enregistrement},
donc un cache par rappel seul retrouve le même enregistrement et sert sa réponse pré-audit (Table~\ref{tab:e2}).

\begin{table}[h]\centering
\begin{tabular}{lrrrr}
\toprule
bras & appels & exact. & \textbf{exact. dérive} & ctx max \\
\midrule
\textbf{Struct} (contrat typé) & \textbf{26} & \textbf{1.00} & \textbf{1.00} & \textbf{290} \\
Invalidant (crochet, sans contrat) & 28 & 1.00 & 1.00 & 290 \\
Naïf & 80 & 1.00 & 1.00 & 290 \\
Long-contexte & 80 & 1.00 & 1.00 & 2062 \\
RAG & 48 & 0.95 & 0.00 & 290 \\
CBR (clé typée, sans re-vérif.) & 24 & 0.95 & 0.00 & 290 \\
Compétence (prose) & 80 & 0.95 & 0.00 & 297 \\
\bottomrule
\end{tabular}
\caption{E2 (simu). Les bras par rappel seul (RAG/CBR/Compétence) servent du périmé ; Invalidant et Struct récupèrent.}
\label{tab:e2}
\end{table}

La lecture est triple. Les \textbf{mémoires par rappel seul servent du périmé} --- le rappel seul ne récupère pas,
l'audit n'entrant jamais dans leur chemin de réutilisation. La \textbf{récupération exige un mécanisme
d'invalidation}, et l'Invalidant comme Struct en ont un (exactitude-dérive 1.00). Ce que le \textbf{contrat
défaisable typé ajoute au rappel codé à la main} : (i) la \emph{sélectivité} --- Struct re-vérifie la post par
entrée (\texttt{satisfies}) et n'évince que les 2 classes \emph{violées} (approve), là où le rappel jette des
classes entières (4 entrées) et paie les re-dérivations supplémentaires (26 vs 28 appels) ; (ii) la
\emph{généralité} --- le même \texttt{assertPost}, agnostique-à-la-prémisse par construction (les expériences
n'exercent qu'un seul type de prémisse, un basculement de drapeau de conformité) ; (iii) la \emph{sûreté de
composition} (§\ref{sec:e3}). L'exécution réelle (Qwen3.6-27B (Q2\_K\_XL, MTP), $N{=}48$) reproduit ceci :
RAG/CBR/Compétence 0.00 ; Invalidant 14 appels / 1.00 ; Struct 13 appels / 2.6\,s / 1.00 / ctx 278 vs
Long-contexte 1304. L'affirmation défendable n'est pas « seul Struct
récupère » mais « la mémoire par rappel seul ne sait pas désapprendre, et un contrat typé déclaratif fournit la
récupération de façon sélective, générale et sûre en composition ».

\subsection{E1 --- amortissement et transfert structurel}
Un domaine de décomposition structurelle (une méthode qui \emph{crée} un sous-graphe avec des identifiants), sur le
moteur complet. Partition : entraînement, apparentés tenus à l'écart, nouveau tenu à l'écart. C'est un contrôle
d'\textbf{existence-et-sûreté sur un petit ensemble} (2 apparentés, 1 nouveau), \textbf{pas un taux} : avec la
transformation relativiser/lier, \emph{toutes} les instances apparentées transfèrent à 0 appel et \textbf{sainement},
le nouveau paie (pas de faux rejeu), totaux 3 appels contre 5. L'ablation sans transformation « touche » mais
rejoue le \emph{mauvais espace d'identifiants} --- \textbf{non sain}. \textbf{Seule la vérification de sûreté}
distingue une réutilisation saine d'un rejeu erroné ; le seul nombre d'appels les classe à égalité. (Un \emph{taux}
de transfert sur de nombreuses méthodes est un travail futur.)

\subsection{E3 --- sûreté de composition}\label{sec:e3}
En composant des paires sur leurs seuls contrats typés (boîte fermée) et en comparant au résultat moteur (moteur
complet), la décision \textbf{coïncide avec la réalité sur chaque paire évaluée, sans faux-admis} --- le
vérificateur n'accepte jamais à tort ; les paires sous-déterminées ou hors-fragment \emph{escaladent}. C'est montré
sur un petit ensemble construit (3 paires sain/non-sain/escalade ; une 4\textsuperscript{e} ajoute l'oracle) : une
\textbf{démonstration d'existence}, pas un taux de population. Chaque barrière est porteuse sur un exemple :
retirer la complétude de cadre manque une écriture non déclarée ; retirer l'étiquette d'effet admet un effet
externe non vérifié ; retirer la détection de cycle admet un cycle couplé rétractable. La décision ne lit que
l'empreinte partagée ; le vérificateur (entailment par domaines abstraits, sain-mais-incomplet) est l'artefact le
plus développé, plutôt sous-évalué ici. (Un corpus plus grand, non trié à la main, est un travail futur.)

\subsection{P4 --- le plafond de couverture \K{}}
Sur une charge mixte (fraction $p$ entièrement typée ; le reste portant un composant en prose qui prime sur la règle
typée), l'appartenance à \K{} décidée par la \textbf{vraie} barrière de canonicalisation, l'amortissement est un
\textbf{gradient en couverture} (approbation : $0\to19\to44\to69\to94\%$ ; tri : $0\to22\to47\to72\to97\%$).
L'exactitude de Struct est $1.00$ à chaque couverture --- la fraction non typée est un \emph{coût} de micro-tâche,
jamais une \emph{falaise} de sûreté. Une variante gloutonne mémoïsant les enregistrements en prose sur leur clé
typée chute à une exactitude égale à la fraction propre : amortir au-delà de \K{} est non sain. Nous sommes
explicites : c'est une \textbf{illustration construite}, pas une mesure réelle --- nous \emph{fixons} $p$ et
\emph{définissons} les enregistrements en prose pour primer, donc « non sain au-delà de \K{} » découle par
construction. Elle établit la \emph{forme} et la sûreté à chaque niveau ; la fraction canonicalisable d'un corpus
réel est dépendante du domaine et non mesurée ici.

\subsection{E5 --- coût de bookkeeping à l'échelle du flux}
Un contrôle de coût de bookkeeping, pas une affirmation sur le passage à l'échelle de la partie difficile : sur un
espace typé de 200 classes avec un audit, quand la longueur du flux $N$ croît de $1{,}320\to20{,}320$
(l'ensemble des classes est fixe ; aucune méthode nouvelle, aucun modèle), Table~\ref{tab:e5}.

\begin{table}[h]\centering
\begin{tabular}{rrrrrr}
\toprule
$N$ & appels Struct & appels/$N$ & appels Naïf & bibliothèque & évincés \\
\midrule
1{,}320 & 202 & 0,153 & 1{,}320 & 200 & 2 \\
5{,}320 & 202 & 0,038 & 5{,}320 & 200 & 2 \\
20{,}320 & 202 & 0,010 & 20{,}320 & 200 & 2 \\
\bottomrule
\end{tabular}
\caption{E5. Les appels de Struct restent constants ; bibliothèque bornée par l'ensemble (fixe) des classes ; éviction sélective.}
\label{tab:e5}
\end{table}

Les appels de Struct restent constants ($\to$ appels/$N\to0$ tandis que Naïf reste à un) ; la bibliothèque est
bornée par le nombre de classes (trivialement) ; une dérive ne rétracte que les classes invalidées
($O(\text{invalidé})$ : 2 sur 200). Coûts par opération : canonicalisation de quelques $\mu$s/appel
(dépend de l'environnement), une passe d'éviction $\approx0{,}5$\,ms. Le contenu honnête est étroit : le bookkeeping
typé n'est pas le goulet. Il ne teste \textbf{pas} l'échelle dans la dimension qui compte (bibliothèque croissante
de méthodes \emph{distinctes}, corpus réel, modèle réel sur tous les bras) --- travail futur.

\subsection{E6 --- tête-à-tête face aux systèmes de mémoire d'agents nommés}\label{sec:e6}
Les références du Dispositif sont génériques (RAG/CBR/Compétence) ; les systèmes attendus en comparaison sont les
systèmes \emph{nommés}. Nous ajoutons une ré-implémentation minimale fidèle de MemGPT/Letta~\cite{packer2023}
(mémoire à étages auto-éditée), Reflexion~\cite{shinn2023} (réflexion verbale pilotée par l'échec, sans mémo) et
GraphRAG~\cite{edge2024} (index de résumés de communautés hors-ligne) derrière la même interface, chacun dans sa
configuration \emph{la plus favorable} avec une ablation appariée (le contrôle négatif). Stub déterministe, $N=78$,
deux classes auditées (Table~\ref{tab:e6}).

\begin{table}[h]\centering
\begin{tabular}{lrrrr}
\toprule
bras & appels & exact. & exact.-dérive & ctx max \\
\midrule
Naïf & 78 & 1,00 & 1,00 & 290 \\
MemGPT (audit paginé en core) & 23 & 1,00 & 1,00 & 320 \\
\quad $-$pagination (ablation) & 18 & 0,92 & 0,00 & 258 \\
Reflexion (signal d'échec différé) & 82 & 1,00 & 1,00 & 332 \\
\quad $-$signal (ablation) & 78 & 0,92 & 0,00 & 258 \\
GraphRAG (index hors-ligne) & 87 & 0,92 & 0,00 & 336 \\
\quad $+$ré-index (ablation) & 89 & 1,00 & 1,00 & 336 \\
\textbf{Struct} & \textbf{20} & \textbf{1,00} & \textbf{1,00} & \textbf{290} \\
\bottomrule
\end{tabular}
\caption{E6. Dans sa configuration la plus favorable chaque système nommé peut récupérer ; seul Struct récupère à la dérive tout en restant non-dominé sur (appels, exact.-dérive, contexte).}
\label{tab:e6}
\end{table}

La lecture honnête n'est \emph{pas} « seul Struct récupère » : dans sa configuration la plus favorable, chaque
système nommé peut récupérer à la dérive. Le propos est que Struct récupère au moindre coût sur les trois axes à la
fois --- c'est le seul bras qui récupère à la dérive tout en restant non-dominé (aucun bras ne l'égale-ou-bat sur
les appels ET le contexte par appel tout en récupérant ; les bras moins chers sont sur la frontière de coût mais
échouent à la dérive). Chaque système nommé paie une taxe propre : MemGPT ne
récupère qu'en paginant l'audit fait-surface en mémoire core (un tour modèle), puis \emph{grossièrement} (un
drapeau de classe entière re-décidant même les membres non concernés) avec un bloc core dans chaque invite ;
Reflexion n'a \emph{aucun} mémo adressé par contenu, donc émet un appel modèle sur \emph{chaque} enregistrement
(appels $\approx N$) et ne récupère que réactivement après un échec observé ; l'index hors-ligne de GraphRAG est
aveugle à l'audit silencieux (périmé), ne récupérant que par un \emph{ré-résumé par lot} (grossier, hors-bande),
jamais un défaiseur par entrée. Les systèmes nommés n'ont pas été conçus pour amortir des décisions ponctuelles
typées récurrentes ; cette charge les éprouve hors de leur domaine de conception, de sorte que la « taxe » est le
coût d'un outil inadapté, non une infériorité à leur propre tâche. Les ablations confirment que chaque mécanisme
est porteur (éteint $\to$ exact.-dérive 0,00). En isolation, la taxe de récupération de Struct est la re-dérivation
des seules entrées \emph{violées} (sélectif) et la ré-assertion du contrat est intra-moteur (\textbf{zéro} appel
modèle). Un essai en direct (Qwen3.6-27B (Q2\_K\_XL, MTP), $N=32$) reproduit l'ordre : Struct 9 appels / 1,9\,s,
MemGPT 11 / 2,3\,s, Reflexion 34 / 7,1\,s, GraphRAG 36 / 7,7\,s (périmé). Ce sont des ré-implémentations minimales
fidèles, pas les systèmes complets.

\subsection{E7 --- composition à la dérive}
E6 mesure une \emph{seule} méthode. Le différenciateur d'une \emph{bibliothèque} est la composition : lorsque la
prémisse d'une méthode amont tombe, la récupération se propage-t-elle dans la chaîne ? Nous étendons la charge à une
chaîne à deux maillons --- \texttt{decide $\to$ disburse}, où \texttt{disbursement = disbursed iff decision ==
approve} (le maillon 2 lit le fait-résultat du maillon 1) --- et ré-exécutons le tête-à-tête en notant les
\textbf{deux} maillons. Le même audit exogène se propage désormais en cascade : un cas audité à score élevé doit
faire basculer \texttt{decision} approve$\to$reject \textbf{et} \texttt{disbursement} disbursed$\to$held.
Simulateur déterministe ($N{=}78$, deux classes auditées) et une exécution réelle (Qwen3.6-27B (Q2\_K\_XL, MTP),
$N{=}48$), Table~\ref{tab:e7}.

\begin{table}[h]\centering
\begin{tabular}{lccc}
\toprule
bras & appels (simu / réel) & exact.-dérive maillon 1 & exact.-dérive maillon 2 \\
\midrule
Naïf & 156 / 96 & 1,00 & 1,00 \\
CBR ($=$ Struct $-$ contrat) & 36 / 24 & \textbf{0,00} & \textbf{0,00} \\
MemGPT (le plus favorable) & 45 / 41 & 1,00 & 1,00 \\
Reflexion (le plus favorable) & 160 / 108 & 1,00 & 1,00 \\
GraphRAG $+$ ré-index & 178 / 116 & 1,00 & 1,00 \\
\textbf{Struct} & \textbf{38 / 26} & \textbf{1,00} & \textbf{1,00} \\
Struct (moteur complet) & 38 / 27 & 1,00 & 1,00 \\
\bottomrule
\end{tabular}
\caption{E7 (simu / réel). Chaîne à deux maillons \texttt{decide $\to$ disburse} : la péremption par rappel seul se cumule aux deux maillons ; Struct récupère les deux sélectivement et est l'unique point Pareto-optimal parmi les bras corrects-à-la-dérive sur (appels, exact.-dérive maillon 1, exact.-dérive maillon 2, contexte).}
\label{tab:e7}
\end{table}

Trois choses sont mesurées. (i) \textbf{La péremption se cumule.} Tout bras par rappel seul ou aveugle (CBR, et
l'ablation de chaque système nommé) est faux au maillon 1 \emph{et} au maillon 2 : une réponse amont périmée
empoisonne le maillon aval qui la lit. L'exactitude à la dérive est mesurée sur l'oracle déterministe ; l'exécution
réelle confirme l'ordre des appels et du temps (le modèle réel est imparfait à petit $N$, nous ne lisons donc pas
la dérive par enregistrement sur lui). (ii) \textbf{La taxe de
récupération se multiplie le long de la chaîne.} Reflexion, sans mémo, paie un appel par enregistrement \emph{par
maillon} (appels $\approx 2N$) ; GraphRAG ré-indexe les communautés concernées à \emph{chaque} maillon ; MemGPT
paie la pagination $+$ un re-décodage grossier $+$ un contexte plus grand aux \emph{deux} maillons. (iii)
\textbf{La cascade de Struct est sélective.} Une prémisse tombée dé-cast la croyance amont et le changement se
propage à l'aval --- récupérant les deux maillons tout en re-dérivant \emph{seulement} l'entrée amont violée (la
re-dérivation aval est élidée car sa clé de cache est l'ensemble de lecture réel de disburse
\texttt{\{kind, region, decision\}}). Struct est l'unique point Pareto-optimal parmi les bras corrects-à-la-dérive
sur (appels $\times$ exact.-dérive maillon 1 $\times$ exact.-dérive maillon 2 $\times$ contexte par appel), en simu
comme en réel ; la réalisation sur
moteur complet (quatre concepts \emph{ensure}-gated au-dessus du cache de dérivation) la reproduit ($38{=}38$ en
simu ; $26\approx27$ en réel, dans le non-déterminisme du modèle). C'est la capacité qui manque structurellement
aux mémoires de surface nommées : une croyance qui dépend d'une prémisse qui vient de tomber, désapprise
\emph{à travers la composition}.

\textbf{Le coin s'élargit avec la profondeur de chaîne.} En généralisant la chaîne à $L$ maillons (chacun positif
ssi le précédent l'est), la péremption et le coût croissent tous deux tandis que la récupération de Struct, non :
la \emph{profondeur de cumul} d'un cache par rappel seul (maillons faux sur une classe dérivée) vaut
exactement $L$ ; Naïf et Reflexion paient $O(L{\cdot}N)$ appels ; mais la \textbf{taxe de dérive} de la
récupération de Struct est en $O(1)$ en $L$ --- la cascade ne re-dérive que l'entrée amont violée, chaque
re-dérivation aval réutilisant l'entrée d'ensemble-de-lecture d'un frère. Ainsi, plus la chaîne de méthodes
apprises est profonde, plus l'avantage de Struct est grand sur l'écart d'exactitude (cumul $\propto L$) et
sur l'efficacité de récupération ($O(1)$ contre une reconstruction grossière en $O(L)$). (La récupération en $O(1)$
tient lorsque chaque re-dérivation aval retombe sur une clé de cache qu'un frère a déjà remplie --- vrai ici : la
branche négative est partagée et les frères à faible score la pré-remplissent ; sous un fort éventail aval, la taxe
de récupération croît vers $O(L)$. Le cumul $\propto L$ découle par construction du fait de définir le maillon $i$
comme une fonction du maillon $i-1$, et est mesuré sur le simulateur déterministe --- le balayage en $L$ en réel est
confondu par le bruit, \S6.)

\textbf{Sur le vrai exécuteur durable.} Ce qui précède est la vue-croyance (le graphe à base de règles $+$ la
rétractation JTMS). La même chaîne compilée en un réseau de workflow et exécutée comme un flot de jetons sur un
magasin de points de reprise adressé par contenu et reprenable après crash reproduit le résultat sur la couche
d'\emph{exécution} : la chaîne amortit et la dérive se propage en cascade à travers les deux maillons
(Struct-sur-exécuteur 24 appels, dérive $1{,}00/1{,}00$ ; la clé sans-prémisse comme le cache composé plat y
cumulent la péremption), la bibliothèque composée chaude \textbf{se rejoue à travers un redémarrage de processus à
zéro appel modèle}, et une chaîne coupée en plein vol \textbf{reprend sans travail perdu ni dupliqué}. La défaisance
en vue-croyance et l'exécution durable sont complémentaires : le contrat rétracte la croyance ; l'exécuteur rend le
recalcul durable et sélectif.

\subsection{E8 --- révision de bibliothèque sous dérive récurrente}\label{sec:e8}
E2/E6/E7 mesurent le chemin \emph{rétracter $\to$ re-dériver} : une post violée évince la croyance en cache, qui est
ensuite re-dérivée. L'autre moitié de la boucle --- \emph{réviser} : spécialiser la précondition de la méthode sur
blâme, de sorte que la bibliothèque soit améliorée, non simplement ré-évaluée --- est évaluée ici. Nous exécutons
$K{=}5$ épisodes de dérive récurrente (les mêmes classes reviennent à chaque épisode après qu'une précondition trop
générale a été exposée) et comparons l'éviction-seule à la révision via le \texttt{reviseOnBlame} du moteur
(Table~\ref{tab:e8}).

\begin{table}[h]\centering
\begin{tabular}{lcc}
\toprule
sur $K{=}5$ épisodes & Éviction-seule & Révision (\texttt{reviseOnBlame}) \\
\midrule
blâmes (par épisode $\to$ cumulatif) & 1,1,1,1,1 $\to$ 5 & 1,0,0,0,0 $\to$ 1 \\
re-dérivations (appels modèle) & 4,1,1,1,1 $\to$ 8 & 4,1,0,0,0 $\to$ 5 \\
taux de faux-admis & 0,25 à chaque épisode & 0,25 $\to$ 0,00 après e1 \\
\bottomrule
\end{tabular}
\caption{E8. L'éviction-seule re-blâme et re-dérive la classe périmée à chaque épisode (coût $\propto K$) ; la révision spécialise la précondition une fois puis se stabilise (0 re-blâme, faux-admis $\to$ 0).}
\label{tab:e8}
\end{table}

L'éviction-seule conserve la règle trop générale, donc elle ré-admet, re-blâme et re-dérive la même classe périmée
à \emph{chaque} épisode (coût $\propto K$). La révision spécialise la précondition une seule fois --- en la
restreignant avec l'atome discriminant du contre-exemple --- puis se stabilise : plus aucun blâme, faux-admis
$\to$ 0. La révision est \emph{chirurgicale} : la précondition spécialisée exclut la classe fautive tout en
admettant encore son frère (elle restreint, elle ne supprime pas la méthode). L'effet tient pour deux types de
prémisse --- un basculement de conformité catégoriel (\texttt{\$compliant != false}) et une porte numérique
(\texttt{\$score != 680}). \textbf{Limite honnête :} \texttt{reviseOnBlame} spécialise par exclusion ponctuelle de
contre-exemple, non par resserrement de borne ; une dérive numérique couvrant $D$ valeurs distinctes converge en
$D$ blâmes ponctuels (bornés, par valeur) plutôt qu'en un seul déplacement de borne --- toujours catégoriquement
mieux que la récurrence par épisode de l'éviction-seule. C'est l'étape qui distingue le contrat de l'invalidation
de cache (\S5) : le blâme alimente la \emph{révision de bibliothèque}, non le seul recalcul de valeurs.

\section{Travaux apparentés}
\textbf{Recherche et mémoire de cas.} RAG~\cite{lewis2020} et CBR~\cite{aamodt1994} rappellent par similarité ; ils ne représentent
pas une \emph{prémisse devenant invalide} (E2). Les bibliothèques de compétences (Voyager~\cite{wang2023}) stockent
de la \emph{prose} sans prémisse défaisable.

\textbf{Mémoire des agents LLM.} MemGPT/Letta~\cite{packer2023}, Reflexion~\cite{shinn2023}, GraphRAG~\cite{edge2024}
gèrent \emph{quoi} garder/rappeler finement, mais par pertinence/récence/similarité ; aucun ne représente une
prémisse typée dont la falsification rétracte une réutilisation. Ils sont complémentaires : un contrat défaisable
pourrait s'y placer comme couche de rétractation. Nous menons ce tête-à-tête en E6 (\S\ref{sec:e6}) : chacun, dans
sa configuration la plus favorable, peut récupérer, mais seul le contrat défaisable le fait au moindre coût sur
(appels $\times$ exactitude $\times$ contexte) simultanément.

\textbf{Apprentissage de bibliothèque / EBL.} DreamCoder~\cite{ellis2021} et Stitch~\cite{bowers2023} font croître une
bibliothèque par abstraction~\cite{plotkin1970} ; l'EBG~\cite{mitchell1986} spécialise depuis une seule preuve ;
aucun n'attache de contrat d'exécution défaisable.

\textbf{Contrats, blâme, vérification graduelle.}~\cite{findler2002,bader2018} sont la lignée de notre discipline,
appliquée aux méthodes \emph{apprises}.

\textbf{Composition.} Non-terminal HRG à deux faces~\cite{habel1992,drewes1997,courcelle1990} à sélection
HTN~\cite{erol1994} ; plan HTN récursif indécidable~\cite{erol1996} ; composition saine via empreintes de
séparation~\cite{ohearn2001,reynolds2002}, problème du cadre~\cite{mccarthy1969} (E3).

\textbf{Révision de théorie/croyances (le voisin le plus proche).} \texttt{reviseOnBlame} relève de la révision de
théorie d'une base de règles \emph{apprise} : EITHER~\cite{ourston1994} et FORTE~\cite{richards1995} révisent des
théories Horn sur contre-exemples ; la révision d'un ensemble de croyances est l'AGM~\cite{agm1985}, et le moniteur
d'exécution s'inscrit dans la tradition du raisonnement défaisable / non-monotone~\cite{reiter1980}. Notre apport
est opérationnel : attacher la révision à une bibliothèque de méthodes \emph{typée, composable, canonicalisable} ;
nous mesurons cette étape de révision sous dérive récurrente en E8 (\S\ref{sec:e8}).

\textbf{Maintien de la vérité.} Un JTMS~\cite{doyle1979,dekleer1986}.

\textbf{Exécution durable / IVM.} Un cas est un marquage 1-sûr sur un réseau de workflow~\cite{vanderaalst1998} ;
l'exécuteur durable est un artefact connu~\cite{aws2022,dbos2022}. DBSP~\cite{budiu2023}/Materialize maintiennent
des \emph{valeurs} ; nous ne revendiquons pas la nouveauté pour \emph{récupérer} sur dérive --- la différence est
que notre invalidation est \emph{dérivée d'un contrat déclaratif} (re-vérifier la post ; blâmer ; spécialiser). Le
rejeu chaud et la reconstruction du seul maillon affecté (\S4.8) s'apparentent au calcul fonctionnel
auto-ajustable~\cite{acar2002} et à la reconstruction incrémentale des systèmes de build~\cite{mokhov2018}.

\section{Menaces à la validité}
\textbf{Simulateur vs réel.} Le simulateur retire l'erreur du modèle pour isoler le \emph{mécanisme} ; l'E2 réel
confirme l'ordre. Le simulateur ne prédit pas l'exactitude absolue en réel ; exécuter plus de bras en réel est un
travail futur. L'exactitude à la dérive est mesurée sur l'oracle déterministe ; les exécutions réelles confirment
l'ordre des appels et du temps (le modèle réel est imparfait à petit $N$, nous ne lisons donc pas la dérive par
enregistrement sur lui).

\textbf{\K{}-couverture paramétrée.} P4 \emph{fixe} la fraction typée ; les affirmations non circulaires sont la
\emph{forme} (un gradient), l'universalité de la sûreté, et la non-sûreté de l'amortissement glouton. La fraction
absolue d'une charge de production est dépendante du domaine.

\textbf{Force des références.} RAG/CBR/Compétence sont nos implémentations ; le bras Invalidant isole « possède un
crochet » de « possède un contrat typé », d'où l'affirmation précise. Une RAG/CBR à invalidation-événementielle
ajustée est essentiellement l'Invalidant. Les systèmes \emph{nommés} sont évalués en E6 (\S\ref{sec:e6}) : des
ré-implémentations minimales fidèles de chaque mécanisme (pas d'édition mémoire réellement pilotée par le LLM, ni
modèle de récompense appris, ni communautés Leiden), chacune dans sa configuration la plus favorable et
\emph{charitable} --- des bornes supérieures du comportement à la dérive, non des hommes de paille. De façon
symétrique, notre STRUCT n'est pas un stub : une réalisation sur le vrai moteur (rétractation JTMS \emph{ensure}-gated
au-dessus du cache de dérivation) reproduit son coin en stub et en réel (9 appels sur le $N=32$ réel) et la
bibliothèque chaude survit à un redémarrage à zéro appel --- E6 confronte les ré-implémentations des systèmes nommés
au \emph{vrai} moteur, pas stub-contre-stub.

\textbf{Nouveauté / positionnement.} Aucun mécanisme n'est nouveau ; le travail est une \emph{composition}. Nous le
positionnons comme une unification opérationnelle, non un nouvel algorithme.

\textbf{Échelle et étendue.} Mécanismes à $N$ modeste sur deux domaines synthétiques ; E5 est un coût de
bookkeeping. Le tête-à-tête face aux mémoires d'agents modernes est désormais E6 (\S\ref{sec:e6}, ré-implémentations
minimales fidèles) ; un corpus réel sur un exécuteur durable et une comparaison face aux systèmes \emph{déployés}
restent des travaux futurs.

\textbf{Sûreté éventuelle, non statique.} L'applicabilité d'une méthode apprise est indécidable en général (Rice) ;
la garantie est une sûreté éventuelle via un moniteur d'exécution porteur au-dessus d'une barrière saine-incomplète.

\textbf{Moteur unique / auteur unique.} Une reproduction indépendante renforcerait les affirmations structurelles.

\section{Conclusion}
La mémoire par rappel seul ne sait pas désapprendre ; les bibliothèques statiques sont saines mais figées. Une
bibliothèque apprise de méthodes typées dotée d'un \textbf{contrat d'exécution défaisable} obtient les deux : elle
amortit et compose sur la structure typée, et lorsqu'une prémisse dérive elle \emph{rétracte avec blâme et révise}
--- récupérant l'exactitude là où la recherche, le CBR et les bibliothèques de compétences servent du périmé. Nos
expériences l'isolent : le rappel seul ne récupère pas ; un crochet d'invalidation le peut ; un contrat typé
\emph{déclaratif} fournit cette récupération de façon sélective, générale et sûre-en-composition, à contexte borné,
sur un moteur réel et confirmé en réel --- borné par une fraction canonicalisable elle-même frontière de sûreté.
Chaque mécanisme relève de l'état de l'art ; la contribution est leur composition en une représentation typée
unique où amortissement, vérification de composition et désapprentissage coïncident. C'est, délibérément, une
synthèse d'ingénierie avec une propriété émergente testable --- un désapprentissage principiel et sélectif ---
plutôt qu'un nouvel algorithme.

\paragraph{Code \& reproductibilité.} Le moteur et un artefact d'expérience autonome sont publics sur
\texttt{github.com/9pings/skynet-graph} --- \texttt{artifact/paper-dll/} : mécanismes de base
(\texttt{workload.js}, \texttt{arms.js}, \texttt{harness.js}, \texttt{e1-transfer.js}, \texttt{e3-compose.js},
\texttt{p4-coverage.js}, \texttt{scale.js}, \texttt{measure-e2-live.js}, \texttt{F6-transfer.js}), le tête-à-tête
E6 (\texttt{named-arms.js}, \texttt{struct-real.js}, \texttt{measure-named-h2h.js}), la suite E7
composition-à-la-dérive (\texttt{composed-*.js}, \texttt{struct-real-composed.js}, \texttt{durable-composed.js},
\texttt{chain-depth.js}) et la révision E8 (\texttt{revise.js}), avec la suite déterministe
\texttt{tests/integration/paper-*.test.js} (\texttt{npm test}). Les exécutions en réel utilisent un endpoint local
compatible OpenAI servant Qwen3.6-27B (Q2\_K\_XL, MTP). Sous licence AGPL-3.0-or-later.

\begin{thebibliography}{99}
\bibitem{aamodt1994} A.~Aamodt, E.~Plaza. Case-Based Reasoning: Foundational Issues, Methodological Variations, and System Approaches. \emph{AI Communications} 7(1):39--59, 1994.
\bibitem{acar2002} U.~A.~Acar, G.~E.~Blelloch, R.~Harper. Adaptive Functional Programming. \emph{POPL} 2002, 247--259.
\bibitem{agm1985} C.~E. Alchourr\'on, P.~G\"ardenfors, D.~Makinson. On the Logic of Theory Change. \emph{J. Symbolic Logic} 50(2):510--530, 1985.
\bibitem{aws2022} AWS. Step Functions Distributed Map. 2022.
\bibitem{bader2018} J.~Bader, J.~Aldrich, \'E.~Tanter. Gradual Program Verification. \emph{VMCAI} 2018, LNCS 10747:25--46.
\bibitem{bourtoule2021} L.~Bourtoule, V.~Chandrasekaran, C.~A.~Choquette-Choo, H.~Jia, A.~Travers, B.~Zhang, D.~Lie, N.~Papernot. Machine Unlearning. \emph{IEEE S\&P} 2021, 141--159.
\bibitem{bowers2023} M.~Bowers et al. Top-Down Synthesis for Library Learning. \emph{POPL} 2023; PACMPL 7(POPL).
\bibitem{budiu2023} M.~Budiu et al. DBSP: Automatic Incremental View Maintenance for Rich Query Languages. \emph{PVLDB} 16(7):1601--1614, 2023.
\bibitem{courcelle1990} B.~Courcelle. The Monadic Second-Order Logic of Graphs I. \emph{Inf. Comput.} 85(1):12--75, 1990.
\bibitem{dekleer1986} J.~de~Kleer. An Assumption-based TMS. \emph{Artificial Intelligence} 28(2):127--162, 1986.
\bibitem{doyle1979} J.~Doyle. A Truth Maintenance System. \emph{Artificial Intelligence} 12(3):231--272, 1979.
\bibitem{drewes1997} F.~Drewes, H.-J.~Kreowski, A.~Habel. Hyperedge Replacement Graph Grammars. In \emph{Handbook of Graph Grammars} Vol.~1, World Scientific, 95--162, 1997.
\bibitem{edge2024} D.~Edge et al. From Local to Global: A Graph RAG Approach to Query-Focused Summarization. arXiv:2404.16130, 2024.
\bibitem{ellis2021} K.~Ellis et al. DreamCoder. \emph{PLDI} 2021.
\bibitem{erol1994} K.~Erol, J.~Hendler, D.~S.~Nau. UMCP: A Sound and Complete Procedure for HTN Planning. \emph{AIPS} 1994.
\bibitem{erol1996} K.~Erol, J.~Hendler, D.~S.~Nau. Complexity Results for HTN Planning. \emph{Ann. Math. AI} 18:69--93, 1996.
\bibitem{findler2002} R.~B.~Findler, M.~Felleisen. Contracts for Higher-Order Functions. \emph{ICFP} 2002, 48--59.
\bibitem{habel1992} A.~Habel. Hyperedge Replacement: Grammars and Languages. LNCS 643, Springer, 1992.
\bibitem{lewis2020} P.~Lewis et al. Retrieval-Augmented Generation for Knowledge-Intensive NLP Tasks. \emph{NeurIPS} 2020.
\bibitem{mccarthy1969} J.~McCarthy, P.~J.~Hayes. Some Philosophical Problems from the Standpoint of AI. \emph{Machine Intelligence} 4, 1969.
\bibitem{mitchell1986} T.~M.~Mitchell, R.~M.~Keller, S.~T.~Kedar-Cabelli. Explanation-Based Generalization: A Unifying View. \emph{Machine Learning} 1(1):47--80, 1986.
\bibitem{mokhov2018} A.~Mokhov, N.~Mitchell, S.~Peyton Jones. Build Systems à la Carte. \emph{Proc. ACM Program. Lang.} 2(ICFP), Art.~79, 2018.
\bibitem{ohearn2001} P.~W.~O'Hearn, J.~C.~Reynolds, H.~Yang. Local Reasoning about Programs that Alter Data Structures. \emph{CSL} 2001, LNCS 2142:1--19.
\bibitem{ourston1994} D.~Ourston, R.~J.~Mooney. Theory Refinement Combining Analytical and Empirical Methods. \emph{Artificial Intelligence} 66(2):273--309, 1994.
\bibitem{packer2023} C.~Packer et al. MemGPT: Towards LLMs as Operating Systems. arXiv:2310.08560, 2023.
\bibitem{plotkin1970} G.~D.~Plotkin. A Note on Inductive Generalization. \emph{Machine Intelligence} 5:153--163, 1970.
\bibitem{reiter1980} R.~Reiter. A Logic for Default Reasoning. \emph{Artificial Intelligence} 13(1--2):81--132, 1980.
\bibitem{reynolds2002} J.~C.~Reynolds. Separation Logic. \emph{LICS} 2002, 55--74.
\bibitem{richards1995} B.~L.~Richards, R.~J.~Mooney. Automated Refinement of First-Order Horn-Clause Domain Theories. \emph{Machine Learning} 19(2):95--131, 1995.
\bibitem{shinn2023} N.~Shinn et al. Reflexion: Language Agents with Verbal Reinforcement Learning. \emph{NeurIPS} 2023.
\bibitem{dbos2022} A.~Skiadopoulos et al. DBOS: A DBMS-oriented Operating System. \emph{PVLDB} 15(1):21--30, 2021.
\bibitem{vanderaalst1998} W.~M.~P.~van~der~Aalst. The Application of Petri Nets to Workflow Management. \emph{J. Circuits, Systems and Computers} 8(1):21--66, 1998.
\bibitem{wang2023} G.~Wang et al. Voyager: An Open-Ended Embodied Agent with Large Language Models. arXiv:2305.16291, 2023.
\end{thebibliography}

\end{document}
